% ===== PRÉAMBULE CANONIQUE — à coller VERBATIM en tête de CHAQUE .tex =====
\documentclass[11pt,a4paper]{article}
\usepackage[utf8]{inputenc}\usepackage[T1]{fontenc}\usepackage{lmodern}
\usepackage[french]{babel}
\usepackage{amsmath,amssymb,mathtools}\usepackage{icomma}
\usepackage{siunitx}\sisetup{locale=FR}  % \qty{}{}, \si{}, \ang{}
\usepackage{xcolor}\usepackage{colortbl}
\usepackage{tikz}\usepackage{pgfplots}\pgfplotsset{compat=1.18}
\usepackage[siunitx,europeanresistors]{circuitikz} % circuits (élec.)
\usepackage[most]{tcolorbox}\usepackage{enumitem}\usepackage{array,booktabs}
\usepackage[a4paper,margin=2.2cm]{geometry}
\usetikzlibrary{arrows.meta,calc,angles,quotes,positioning,patterns,%
  backgrounds,shapes.geometric,decorations.pathmorphing,decorations.markings,intersections}
\definecolor{primary}{RGB}{222,49,99}\definecolor{primarydark}{RGB}{150,22,55}
\definecolor{defcol}{RGB}{0,114,178}\definecolor{methcol}{RGB}{52,92,148}
\definecolor{excol}{RGB}{0,135,90}\definecolor{attncol}{RGB}{200,40,55}
\tcbset{boxbase/.style={enhanced,breakable,boxrule=0.9pt,arc=2.5pt,
  left=3mm,right=3mm,top=2.3mm,bottom=2.3mm,fonttitle=\bfseries,coltitle=white}}
% --- boîtes TUTORIEL (noms EXACTS, ne pas renommer) ---
\newtcolorbox{cle}[1][]{boxbase,colback=defcol!6,colframe=defcol,
  title={\textsc{À retenir}\ifx\relax#1\relax\relax\else\textmd{ : \itshape #1}\fi}}
\newtcolorbox{methode}[1][]{boxbase,colback=methcol!7,colframe=methcol,sharp corners,
  title={\textsc{Méthode}\ifx\relax#1\relax\relax\else\textmd{ --- \itshape #1}\fi}}
\newtcolorbox{exemplebox}[1][]{boxbase,colback=excol!5,colframe=excol,
  title={\textsc{Exemple}\ifx\relax#1\relax\relax\else\textmd{ : \itshape #1}\fi}}
\newtcolorbox{piege}[1][]{boxbase,colback=attncol!6,colframe=attncol,
  title={\textsc{Piège}\ifx\relax#1\relax\relax\else\textmd{ : \itshape #1}\fi}}
\newtcolorbox{remarque}[1][]{boxbase,colback=black!4,colframe=black!45,
  title={\textsc{Remarque}\ifx\relax#1\relax\relax\else\textmd{ : \itshape #1}\fi}}
% --- boîtes EXERCICE / CORRIGÉ (noms EXACTS, auto-comptées) ---
\newtcolorbox[auto counter]{exo}[1][]{boxbase,colback=primary!4,colframe=primary,
  title={\textsc{Exercice}~\thetcbcounter\ifx\relax#1\relax\relax\else\textmd{ --- \itshape #1}\fi}}
\newtcolorbox[auto counter]{corrige}[1][]{boxbase,colback=excol!4,colframe=excol,
  title={\textsc{Corrigé}~\thetcbcounter\ifx\relax#1\relax\relax\else\textmd{ --- \itshape #1}\fi}}
\newcommand{\vv}[1]{\vec{#1}}\newcommand{\ds}{\displaystyle}
\newcommand{\R}{\mathbb{R}}\newcommand{\N}{\mathbb{N}}\newcommand{\Z}{\mathbb{Z}}
% (un \usepackage{fancyhdr}+en-tête est possible ; il est ignoré côté web)
% =========================================================================
\begin{document}

\begin{center}
{\large\bfseries\color{excol} Thème 3 — Corrigé Difficile}\\[-2pt]
{\color{excol}\rule{5cm}{1pt}}
\end{center}

\begin{corrige}[géométrie d'un mode]
4 fuseaux $\Rightarrow n=4$, donc $\lambda_n=2L/n$ avec $n=4$.
\begin{enumerate}[label=\alph*)]
  \item $\lambda_4=\dfrac{2L}{4}=\dfrac{2\times 2}{4}=\SI{1}{\metre}.$
  \item Nœud--nœud $=\dfrac{\lambda_4}{2}=\SI{0.5}{\metre}$ ;
        nœud--ventre voisin $=\dfrac{\lambda_4}{4}=\SI{0.25}{\metre}.$
  \item $n+1=4+1=\textbf{5 nœuds}$ (les deux extrémités comprises).
\end{enumerate}
\end{corrige}

\begin{corrige}[identifier un harmonique sur un schéma]
On compte les fuseaux : il y en a 4 (et 5 nœuds).
\begin{enumerate}[label=\alph*)]
  \item $n=4$ : c'est la 4\textsuperscript{e} harmonique. $\lambda_4=\dfrac{2L}{4}=\dfrac{L}{2}.$
  \item $f_4=n f_1=4\times 50=\SI{200}{\hertz}.$
\end{enumerate}
\end{corrige}

\begin{corrige}[remonter à la célérité]
5 fuseaux $\Rightarrow n=5$. On lit $\lambda$ via $\lambda_n=2L/n$, puis $c=\lambda f$.
\begin{enumerate}[label=\alph*)]
  \item $\lambda_5=\dfrac{2L}{5}=\dfrac{2\times 1}{5}=\SI{0.4}{\metre}.$
  \item $c=\lambda_5\,f=0{,}4\times 250=\SI{100}{\metre\per\second}.$
  \item $f_1=\dfrac{c}{2L}=\dfrac{100}{2}=\SI{50}{\hertz}$
        (vérification : $f_5=5\times 50=\SI{250}{\hertz}$, cohérent).
\end{enumerate}
\end{corrige}

\begin{corrige}[accorder une corde]
La fondamentale $f_1=\dfrac{c}{2L}$ augmente si $L$ diminue \emph{ou} si $c$ augmente.
\begin{enumerate}[label=\alph*)]
  \item $f_1=\dfrac{c}{2L}=\dfrac{143}{2\times 0{,}65}=\dfrac{143}{1{,}3}=\SI{110}{\hertz}.$
  \item $L'=L/2\Rightarrow f_1'=\dfrac{c}{2L'}=\dfrac{c}{L}=\dfrac{143}{0{,}65}=\SI{220}{\hertz}$ :
        la fréquence \textbf{double} (une octave plus haut).
  \item $c'=2c\Rightarrow f_1=\dfrac{2c}{2L}=\dfrac{286}{1{,}3}=\SI{220}{\hertz}$ : elle
        double également. Raccourcir la corde ou la tendre montent tous deux la note.
\end{enumerate}
\end{corrige}

\begin{corrige}[tuyau / corde à extrémité libre]
Le bout libre impose un ventre : seuls les modes impairs $2n-1=1,3,5,\dots$ existent.
\begin{enumerate}[label=\alph*)]
  \item $f_1=\dfrac{(2\times1-1)c}{4L}=\dfrac{c}{4L}=\dfrac{300}{4\times 0{,}75}=\dfrac{300}{3}=\SI{100}{\hertz}.$
  \item $f_1=\SI{100}{\hertz}$, $f_2=3f_1=\SI{300}{\hertz}$, $f_3=5f_1=\SI{500}{\hertz}$
        (multiples \emph{impairs} de $f_1$).
  \item $2f_1$ correspondrait à un mode « pair », qui exigerait un \emph{nœud} à l'extrémité
        libre. Or une extrémité libre est nécessairement un \textbf{ventre} : ce mode est donc
        impossible.
\end{enumerate}
\end{corrige}

\end{document}
